% ============================================================
% discussion.tex  [ENGLISH — RCR submission version]
% Paper: Municipal Public Cleaning Expenditure and Solid Waste
%        Management: Evidence from Peru, 2015–2019
% Author: Dante Barreto Gamarra (UNAS)
% ============================================================

\section{Discussion}
\label{sec:discusion}

\subsection{Effect Magnitudes and Comparative Benchmarks}
\label{sec:benchmarks}

The estimated elasticity of $+0.879$ between municipal expenditure and daily waste collected falls in the upper range of effects documented in the solid waste management cost and efficiency literature. \citet{Bel2010}, in their analysis of Catalan municipalities, document that economies of scale in waste management are exhausted above populations of approximately 20,000 inhabitants, implying decreasing expenditure-output elasticities in larger municipalities. The elasticity of $+0.879$ estimated for Peru is consistent with that framework: municipalities operating below the minimum efficient scale threshold exhibit near-proportional returns. \citet{Kinnaman2006} argues, from a welfare economics perspective, that underprovision of collection services in low-income municipalities reflects budget constraints rather than preferences; the results of this paper support that reading for the Peruvian case: the expenditure-QPRS elasticity close to unity ($+0.879$) indicates that collection capacity is approximately proportional to expenditure, so the coverage gap observed between rich and poor municipalities is primarily attributable to the fiscal gap, not to differences in technical efficiency.

The elasticity of $+0.879$ falls in the mid-to-high range of the comparative evidence from low- and middle-income countries (LMICs), and its position within that range is interpretable in terms of institutional structure.
\citet{Guerrero2013}, in a comparative study of waste management systems in cities across Latin America, Sub-Saharan Africa, and Asia, document that the primary constraint for formal collection in low-fiscal-capacity municipalities is operational financing, not technological availability or the regulatory framework.
In the Sub-Saharan contexts analyzed --- where fiscal decentralization has advanced in parallel with the transfer of solid waste competencies to local governments, as in Ghana, Kenya, and Tanzania --- cities with more stable intergovernmental transfers exhibit collection coverage rates 15 to 30 percentage points higher than comparable cities with volatile financing.
Evidence from the \textit{What a Waste 2.0} report \citep{Kaza2018} corroborates this pattern for 368 cities in low- and lower-middle-income countries: collection coverage is strongly correlated with municipal per capita operating expenditure, with a log-log slope of $0.7$--$0.9$ for the poorest city quartile --- a range that coincides with the Peruvian elasticity of $+0.879$.

The contrast with South Asian contexts, particularly India, illustrates why the Peruvian elasticity is mid-to-high rather than low.
In secondary urban municipalities of Indian states with mixed formal-informal provision, marginal fiscal expenditure faces a substitution effect with the informal sector that depresses the effective elasticity: when the municipality formalizes collection, it displaces pre-existing informal providers, and the net increase in waste collected is smaller than the additional expenditure would suggest \citep{Kaza2018}.
In rural Peru, that substitution is minimal: the smallest district municipalities have no organized informal collection sector, so each additional sol in expenditure translates directly into higher coverage with no displacement effect.
This explains why the Peruvian elasticity exceeds typical South Asian values ($0.40$--$0.60$) and aligns with Sub-Saharan Africa ($0.70$--$0.95$), where fiscal decentralization is more recent and informal provision less organized.
The solid waste management gap among Peruvian municipalities is, in this sense, primarily a budget constraint rather than a technical gap: operational capacity scales with financing because there is no alternative provision to substitute for it.

The average marginal effect on the probability of sanitary landfill use (AME $= +0.022$ per 10\% expenditure increase) is modest at the individual level but cumulatively relevant at the national aggregate. Projected across the 1,874 municipalities in the sample, a sustained 10\% increase in public cleaning expenditure would imply that approximately 41 additional municipalities would transition to formal disposal in a typical year. The AME magnitude is comparable to the effects of direct disposal infrastructure subsidies documented in comparable economies: \citet{Kinnaman2006} reports that controlled landfill subsidy programs in budget-constrained jurisdictions produced formal disposal transitions of 1.5 to 3 percentage points per budget cycle. Peru's AME of $+0.022$, obtained without explicit infrastructure subsidies but only through municipal operating expenditure, is comparable in magnitude and reinforces the hypothesis that financing is the binding constraint for environmental formalization in low-fiscal-capacity municipalities.

% ------------------------------------------------------------------
\subsection{External Validity}
\label{sec:validez_externa}

The results of this paper are generalizable to other fiscal decentralization contexts in Latin America with comparable institutional structures.
\citet{Faguet2004}, in his analysis of the Bolivian decentralization process following the 1994 Law of Popular Participation, documents that the transfer of competencies to local governments increased public investment in basic services --- education, health, sanitation --- more responsively to local needs than prior centralized expenditure.
The mechanism identified by \citet{Faguet2004} is analogous to the one documented here: municipalities with greater fiscal resource access reallocate spending toward services with the largest relative deficits, and the binding constraint for this reallocation is financial, not informational.
The Peruvian case extends that finding to public cleaning services: the near-unit elasticity ($+0.879$) indicates that operational capacity scales almost proportionally with expenditure, implying that the coverage gap between rich and poor municipalities is fundamentally a fiscal gap transferable through intergovernmental transfer policy.

The comparison with low- and middle-income countries outside Latin America suggests the results are not idiosyncratic to the Peruvian context.
\citet{Guerrero2013}, in a comparative analysis of waste management systems in developing-country cities across Latin America, Sub-Saharan Africa, and Asia, document that the primary constraint for formalizing final disposal is municipal operating financing, not technology availability or regulation.
In the Sub-Saharan and South Asian contexts analyzed, cities with more stable fiscal transfers exhibit formal disposal rates 15 to 30 percentage points higher than comparable cities with volatile financing, holding per capita income constant.
The AME of $+0.022$ estimated for Peru falls in the lower end of this interval, consistent with the fact that FONCOMUN transfers are relatively stable in real terms but insufficient in absolute level for the majority of rural district municipalities.
The convergence of evidence across such diverse contexts reinforces the interpretation of financing as the universal binding constraint for environmental formalization in middle- and low-income countries.

% ------------------------------------------------------------------
\subsection{Transmission Mechanisms}
\label{sec:mecanismos}

The results are consistent with two distinct transmission channels operating in parallel but with different time horizons. The direct operational channel links expenditure to collection frequency (FR) and quantity collected (QPRS): financing current inputs --- fuel, labor, fleet maintenance --- determines the immediate operational capacity of the service. This channel explains why FR and QPRS exhibit the largest effects and fastest materialization in the panel: they are the only outcomes that respond to current-period expenditure without requiring an intermediate bureaucratic or investment process.

The institutional channel operates through planning. Expenditure enables the hiring of consultants, the preparation of diagnostics, and the citizen participation processes that are prerequisites for the approval of instruments such as the PIGARS. Planning, in turn, is an administrative prerequisite for accessing final disposal infrastructure financing: under Peru's regulatory framework (Legislative Decree No.\ 1278 and its Regulation, DS No.\ 014-2017-MINAM), municipalities without an approved PIGARS are not eligible for MINAM sanitary landfill financing. This generates a causal chain expenditure $\to$ PI $\to$ RS that the model estimates in reduced form, without imposing mediational structure, but which the coefficient pattern suggests: PI responds to expenditure with the largest relative magnitude among all outcomes ($\hat{\beta} = +0.300$), and RS responds with smaller magnitude but statistically unambiguous significance (AME $= +0.022$). The weaker effect of RS relative to PI is consistent with the hypothesis that planning is a necessary but not sufficient condition for formal infrastructure adoption: larger-scale physical investments and longer time horizons are also required.

The strength of the effect on PI ($\hat{\beta}_{\text{CRE+CF}} = 0.300$ vs.\ $\hat{\beta}_{\text{TWFE}} = 0.039$, a ratio of $7.7\times$) reflects that the institutional channel is the most severely distorted by endogeneity. Municipalities with greater managerial capacity tend simultaneously to execute higher budgets and to hold more planning instruments, so the uncorrected TWFE captures that selection covariance as if it were an expenditure effect. Once endogeneity is corrected, the pure effect of financing on planning is substantially larger, indicating that expenditure is itself an enabler of institutional capacity rather than merely a correlate of it.

% ------------------------------------------------------------------
\subsection{Policy Implications}
\label{sec:politica}

The heterogeneity implicit in the near-unit elasticity ($\hat{\beta}_{\text{QPRS}} = 0.879 < 1$) has a direct distributional consequence for budget allocation: marginal returns to public cleaning expenditure are decreasing, so the infra-marginal reallocation of resources from high-expenditure municipalities toward low-expenditure ones would produce net gains in aggregate waste management. This is an efficiency prescription within the existing transfer system, requiring no increase in total expenditure.

FONCOMUN --- Peru's main intergovernmental transfer, distributed by a formula based on population, poverty, and rurality --- is the most relevant fiscal instrument for implementing such reallocation. Since district municipalities with lower fiscal capacity depend on FONCOMUN for 80--95\% of their revenues, a modification of the formula coefficients to weight the solid waste management deficit (measured by the PI index or the open dumpsite rate) would increase financing precisely in municipalities where the marginal return on expenditure is highest. This proposal is consistent with the results-based budgeting logic introduced by DS No.\ 004-2013-EF on performance budgeting in Peru.

An order-of-magnitude estimate quantifies the fiscal cost of the reform.
In 2015--2019, the median public cleaning expenditure for first-quintile fiscal municipalities was approximately S/\,85,000 annually, compared to S/\,1,200,000 for fifth-quintile municipalities (source: SIAF/MEF).
A 10\% increase in first-quintile expenditure would require S/\,8,500 additional per municipality-year.
With an AME of $+0.022$, this increase would produce a formal disposal transition in 2.2 out of every 100 lower-quintile municipalities.
Extrapolated to the 375 first-quintile municipalities, the increase would generate approximately 8 additional transitions per year at a total fiscal cost of S/\,3.2 million annually --- equivalent to $0.006$\% of the Ministry of Economy and Finance budget in the same period.
For the three lower quintiles combined (where marginal expenditure returns are highest according to the $+0.879$ elasticity), a 10\% increase in cleaning expenditure would require S/\,28.4 million in additional fiscal effort annually, producing approximately 41 transitions to formal disposal --- precisely the number projected nationally in \cref{sec:benchmarks}.
Modifying the FONCOMUN formula coefficients to weight the formal disposal deficit would therefore require an amount representing less than 0.05\% of the total intergovernmental transfer budget, placing it within the range of fiscally neutral reforms implementable without increasing total expenditure.

However, expenditure increases without parallel institutional reform face structural limits. The AME on RS ($+0.022$ per 10\% increase) implies that moving a municipality from open dumping to sanitary landfill requires expenditure increases that exceed the fiscal capacity of most rural district municipalities under the current scheme. Final disposal infrastructure requires capital investment with economies of scale unachievable at the individual district level: joint municipal landfills (mancomunidades), provided for under Law No.\ 29029 (Municipal Association Law) and MINAM guidelines, are the policy solution the data support. The positive effect on PRS (AME $= +0.026$) among municipalities already accessing sanitary landfills indicates that the problem of under-utilization of existing infrastructure also persists, and can be addressed through increases in operating expenditure without new capital investment.

% ------------------------------------------------------------------
\subsection{Limitations}
\label{sec:limitaciones}

The analysis period (2015--2019) predates the COVID-19 pandemic, which discontinuously altered both solid waste generation patterns and municipal budgets in Peru. The estimated effects are not extrapolable to the post-2020 period without verifying that the expenditure-management relationship structure was not modified by that shock. Comparison with RENAMU 2020--2023 data, once available and audited by INEI, will allow assessment of parameter temporal stability.

The CRE+CF estimator controls for unobservable institutional heterogeneity of a permanent nature, but does not purge time-varying confounders without an external instrument; the integration of FONCOMUN as an instrumental variable --- pending the merge with MEF/SIAF data --- would strengthen identification in future work. The available FONCOMUN instrumentation produces a Kleibergen-Paap $F = 2.8$--$6.5$, below the Stock-Yogo 10\% threshold ($F = 16.4$); the 2SLS bias from a weak instrument with positive correlation with expenditure is directed toward the OLS estimator, making the IV estimator more conservative than the CRE+CF reported --- not more aggressive. To quantify sensitivity to residual unobservable selection, we calculate the \citet{Oster2019} plausibility bounds: the $\delta^*$ statistic is negative and greater than 1.3 in absolute value for all three primary outcomes (Appendix~\ref{sec:appendix_oster}), implying that unobservable selection would need to operate in the opposite direction from observable selection --- a condition contrary to the bias direction documented by the CF test --- to reduce the effects to zero.

The management variables come from RENAMU (National Municipal Registry), a statutory declaration without independent physical audit. Data quality depends on the truthfulness of municipal officials in completing INEI forms. The likely direction of measurement bias depends on whether reporting quality is orthogonal to expenditure or correlated with it. If reporting quality is primarily determined by administrative capacity --- which the Mundlak device absorbs as permanent heterogeneity --- the remaining bias would be classical attenuation in the dependent variable, compressing coefficients toward zero (conservative bias). If, alternatively, higher-expenditure municipalities report better indicators in an aspirational manner correlated with transitory expenditure $\tilde{G}_{it}$, the error would generate upward bias in the reported coefficients. The first hypothesis is more plausible given that the RENAMU form is administered by INEI with a standardized protocol, but both possibilities are non-dismissible without a field-data validation study.

External socioeconomic controls --- provincial poverty (ENAHO), urbanization and population (CPV 2017) --- were linearly interpolated between survey years, introducing measurement error into the control regressors. Since these controls are not the variable of interest, the effect on $\hat{\beta}$ estimates depends on the degree of correlation between expenditure and the interpolated controls; the sensitivity analysis in the Appendix suggests that the direction and magnitude of the main coefficients are stable to variations in the controls.

% ------------------------------------------------------------------
\subsection{Research Agenda}
\label{sec:agenda}

The most immediate and highest-value identification extension is the merge of SIAF/MEF expenditure data with FONCOMUN data for the same period. With that base, the FONCOMUN distribution formula --- whose coefficients are set centrally and do not respond to individual municipality decisions --- can be used as an external instrument for $\ln G_{it}$, transforming the CRE+CF estimator into an IV-2SLS within the CRE framework. The IV estimator would recover the LATE for the subgroup of municipalities whose public cleaning expenditure responds to FONCOMUN transfer variations --- predictably those with lower fiscal capacity and greater transfer dependence, which are precisely those with the highest marginal expenditure returns according to the estimated elasticity.

Heterogeneity analysis by municipality size and geographic region constitutes a second priority extension. The aggregate elasticity of $0.879$ conceals an effect distribution that may vary substantially between large urban municipalities (where management problems are of scale and logistics) and small rural municipalities (where the constraint is one of basic operational capacity). A specification with interactions between $\ln G_{it}$ and population quartiles, or between $\ln G_{it}$ and natural region dummies (coast, sierra, jungle), would identify whether decreasing expenditure returns operate homogeneously or whether subpopulations exist with significantly higher elasticities where fiscal intervention would have greater impact.

The temporal extension of the panel to 2020--2023, conditional on availability of the corresponding RENAMU, would enable two additional analyses. First, evaluating whether the COVID-19 shock generated a structural break in the expenditure-management relationship, which has implications for the external validity of the estimated parameters. Second, exploiting variation in municipal fiscal incentive programs (Municipal Management Improvement Incentive Plan) that the MEF modified during that period, generating an additional source of quasi-experimental variation in cleaning expenditure.
